\section{实验原理}
\subsection{局部性原理(Principle of Locality)}
Cache存在的理论基础就是局部性原理:
\begin{itemize}
    \item 时间局部性 (Temporal Locality)：如果一个存储单元被访问，则近期内它很可能再次被访问（例如：循环变量、重复调用的函数）。
    \item 空间局部性 (Spatial Locality)：如果一个存储单元被访问，则其邻近的存储单元近期内也很可能被访问（例如：数组遍历、顺序执行的指令）。
\end{itemize}
Cache 通过存放这些“活跃”数据，减少了 CPU 访问高延迟 DRAM 的次数,大大提高了系统运行效率。

\subsection{Cache结构与地址拆分}
假设Cache的块的大小为$B = 2^b$Byte,共有$S = 2^s$个组,每组有$W$路,则对于一个物理地址Addr,其逻辑地址被划分为三部分:
\[  Addr = \text{[Tag|Index|Offset]}   \]
\begin{itemize}
    \item Offset(块内偏移):$b$位,决定数据在Cache\ Line中的具体字节位置。
    \item Index(组索引):$s$位,用于定位该地址属于哪一个Set。计算公式为:
    \[  Index = (Addr >> b)\&(S-1) \]
    \item Tag(标志位):剩余的高位,用于区分映射到同一个组的不同内存块。
\end{itemize}

\subsection{Cache块映射方式(Address Mapping)}
Cache地址映射方式决定了主存块如何被放置到Cache中,是Cache实现的核心:
\begin{itemize}
    \item 直接映射(Direct Mapping)
    \begin{itemize}
        \item 原理:每个主存块映射到Cache中的唯一固定一行(每组只有一行)。
        \item 计算逻辑:$Cache\text{行号} = \text{(主存块号)} \pmod {\text{Cache总行数}}$
        \item 地址划分:
        \begin{itemize}
            \item Offset:块内偏移,由块大小决定。
            \item Index:索引位:直接对应Cache行号。
            \item Tag:高位剩余部分,用于确定该行存放的是否是目标块。
        \end{itemize}
        \item 特点:硬件电路最简单,但冲突试效率高。如果两个经常访问的块映射到同一行,会发生剧烈抖动。
    \end{itemize}
    \item 全相联映射(Fully Associative Mapping)
    \begin{itemize}
        \item 原理:主存中的任意一块可以放置在 Cache 中的任意一行(因为只有一个组)。
        \item 计算逻辑:没有 Index 字段，只需匹配 Tag。
        \item 地址划分:只分为Tag和Offset。
        \item 特点:冲突概率最低,空间利用率最高。但因为全相联高速缓存电路必须并行地搜索许多相匹配的标记,构造
一个又大又快的相联高速缓存很困难,而且很昂贵,硬件成本极高,功耗大,通常只用作极小容量的存储(如TLB)。
    \end{itemize}
    \item 组相联映射(Set-Associative Mapping)
    \begin{itemize}
        \item 原理:这是前两者的折中。Cache 分为若干个组(Set)，每组包含 $k$ 个路(Way)。主存块映射到固定的组，但可以放在该组内的任意一路。
        \item 计算逻辑:$Set Index = (\text{主存块号}) \pmod { Cache\text{组数}}$
        \item 特点:兼顾了直接映射的低成本搜索和全相联的低冲突概率。实现时先通过组索引计算组号,组中任何一行都可以包含任何映射到这一组的内存块,所以
必须搜索每组中每一行,寻找一个有效的行,其标记与地址中的标记相匹配。
    \end{itemize}
\end{itemize}

\subsection{替换算法(Replacement Algorithms)}
当Cache发生缺失(Miss)时,且计算得到的Cache对应位置(组)已经满时,必须决定换出哪一行。
\begin{itemize}
    \item FIFO(First-In-First-Out,先进先出)
    \begin{itemize}
        \item 原理:总是替换掉最先进入Cache的那一块。
        \item 实现逻辑:为每一行维护一个时间戳或使用一个循环队列,最早调入的块被最先替换掉。
        \item 评价:实现简单,但忽略了局部性原理,因为一个很早被调入的块也可能是会被频繁访问的块,这样一个频繁被访问的块也会被替换掉,因此该算法性能一般。 
    \end{itemize}
    \item LRU(Least Recently Used,最近最少使用)
    \begin{itemize}
        \item 原理:替换掉距离现在最长时间未被访问的那一块。
        \item 实现逻辑:每次命中某行时,更新其计数器或移动到链表首部;发生替换时选择计数器数值最小、反映"最久远"或链表尾部的行。
        \item 评价:LRU具体实现会较为复杂,但能很好地利用局部性原理,算法性能较优。
    \end{itemize}
\end{itemize}

\subsection{Cache写策略(Writing Policy)}
Cache的写操作分为写命中(write hit)和写不命中两种情况:
\begin{itemize}
    \item 写命中:
    \begin{itemize}
        \item 直写(write-through):立即将字$w$的告诉缓存块写回紧接的低一层,实现简单,但每次写都会引起总线流量。
        \item 写回(write-back):尽可能迟地更新,只有当替换算法将这个更新过的块要替换出去时才把它写到紧接着的低一层中。该策略利用了局部性原理,能显著减少总线流量;缺点是显著增加了复杂性。
高速缓存必须为每一个高速缓存行维护一个额外的修改位(dirty bit),表明这个块是否被修改过。
    \end{itemize}
    \item 写不命中:
    \begin{itemize}
        \item 写分配(write-allocate):加载相应的低一层的块到Cache中并更新这个Cache块。该方法利用了写的空间局部性,但缺点是每次写不命中都会导致
一个块从低一层传送到Cachhe。
        \item 非写分配(not-write-allocate):避开Cache,直接将这个字写到低一层中。
    \end{itemize}
    \item 直写Cache通常是非写分配的;写回Cache通常是写分配的。

\end{itemize}

